\section{Full native semi-regular torus parent spaces}
\label{sec:peps_native_semiregular_parent}

Let $V=\mathbb Z/m\mathbb Z\times\mathbb Z/n\mathbb Z$, with
$m,n\ge3$, and join nearest neighbors to form the native simple-graph
torus. Each site has legs ordered as top, right, down, left.
The native arrows point right and down, so top and left are incoming.
Throughout this section one matrix representation $U$ is used on all
edges and one tensor $a$ is repeated at every vertex.
The final dimension and spanning results are the corresponding
uniform native-torus statements of
\cite[Theorems~5.7 and~5.9]{Schuch2010PEPS}.
The restrictions to a uniform representation, a repeated tensor, and
periods at least three are retained; the link-dependent and smaller-period
cases are discussed in
\cite{gap:rmp_peps_quantum_double_g_isometry,gap:rmp_peps_examples_small_torus}.
The oriented graph comparison in
Theorem~\ref{thm:peps_oriented_semiregular_physical_parent}, in contrast,
allows vertex-dependent physical tensors.

\subsection{Native geometry and the local invariant subspace}

\begin{definition}[Native torus tails and edge reversals]%
    \label{def:peps_native_torus_orientation}
    \lean{TNLean.PEPS.torusNativeEdgeTail,
        TNLean.PEPS.torusNativeEdgeFlip,
        TNLean.PEPS.torusNativeEdgeTail_right,
        TNLean.PEPS.torusNativeEdgeTail_up}
    \leanok
    For the horizontal edge from $v=(x,y)$ to $(x+1,y)$ the native
    tail is $v$. For the vertical edge joining $v$ to $(x,y+1)$,
    the native tail is $(x,y+1)$, including wraparound edges.
    If $e=(t,h)$ is its ordered-edge description, define $o_e=1$
    exactly when $t$ is not this native tail.
\end{definition}

\begin{theorem}[Native outgoing incidences and degree]%
    \label{thm:peps_native_torus_incidence}
    \lean{TNLean.PEPS.torusNativeEdgeTail_mem_endpoints,
        TNLean.PEPS.card_torusIncidentEdge,
        TNLean.PEPS.torusNativeEdgeTail_leg_iff,
        TNLean.PEPS.torusNativeEdgeTail_incident_iff,
        TNLean.PEPS.graphNativeTail_torus_iff,
        TNLean.PEPS.graphNativeTail_torus_leg_iff}
    \leanok
    \uses{def:peps_native_torus_orientation,def:peps_oriented_native_tail}
    Every native tail is an endpoint of its edge. Every vertex has
    four distinct incident edges, bijectively indexed by
    $0,1,2,3$ in top, right, down, left order.
    At $v$, the native-tail incidences are precisely indices $1,2$.
    Equivalently, the graph native-tail predicate for $o$ holds exactly
    when the chosen torus tail equals $v$.
\end{theorem}
\begin{proof}\leanok
    The horizontal and vertical edge bijection identifies the selected
    tails with actual endpoints. The four-leg incidence bijection gives
    degree four. Evaluating the tail on its four legs gives $v$ on
    right and down; on top and left it gives a neighboring vertex,
    different from $v$. The definition of $o_e$ then identifies this
    test with the ordered-graph native-tail predicate.
\end{proof}

\begin{definition}[Native incidence coordinates and site tensor]%
    \label{def:peps_native_incident_coordinates}
    \lean{TNLean.PEPS.torusNativeIncidentCoordinatesEquiv,
        TNLean.PEPS.torusNativeIncidentCoordinatesEquiv_apply,
        TNLean.PEPS.torusNativeIncidentSite}
    \leanok
    \uses{thm:peps_native_torus_incidence}
    The bijection $j_v:X^{\operatorname{Inc}(v)}\to X^4$ sends
    $\eta$ to $(\eta_{\mathrm{top}},\eta_{\mathrm{right}},
        \eta_{\mathrm{down}},\eta_{\mathrm{left}})$.
    A native tensor $a:X^4\times\{0,\ldots,d-1\}\to\C$
    defines the graph site $a_v(\eta,s)=a(j_v\eta,s)$.
    Its native linear map is
    $\mathcal A x(s)=\sum_{\xi\in X^4}a(\xi,s)x(\xi)$.
\end{definition}

\begin{theorem}[Native four-leg and oriented graph $G$-injectivity]%
    \label{thm:peps_native_oriented_g_injectivity}
    \lean{TNLean.PEPS.graphOrientedIncidentMatrix_torus,
        TNLean.PEPS.graphOrientedIncidentRepresentation_torus,
        TNLean.PEPS.regularSiteMap_torusNativeIncidentSite_eq_siteMap,
        TNLean.PEPS.IsGInjective.isGInjective_torusNativeIncidentSite,
        TNLean.PEPS.localTensorMap_torusNativeIncidentSite,
        TNLean.PEPS.IsGInjective.isGInjective_torusNativeIncidentTensor}
    \leanok
    \uses{def:peps_native_incident_coordinates,
        def:peps_oriented_incident_representation,def:peps_g_injective}
    Define the native four-leg representation by
    \begin{align}
        \rho_\square(g)=U_g\otimes U_{g^{-1}}^T
        \otimes U_{g^{-1}}^T\otimes U_g.
        \label{eq:peps_native_four_leg}
    \end{align}
    The site coefficient $a(\eta_t,\eta_r,\eta_b,\eta_l;s)$ has the
    following native leg orientation; the transverse physical index is $s$.
    % Formula: a(eta_t,eta_r,eta_b,eta_l;s), with incoming top/left and
    % outgoing right/down. The action on these four indices is exactly
    % rho_square(g)=U_g tensor U_(g^-1)^T tensor U_(g^-1)^T tensor U_g.
    % Source: SCP10 paper_v3.tex:1278--1296, eq:2d-ug-sym.
    % Ink-to-index: 90/0/270/180 degrees are top/right/down/left;
    % the transverse policy leg is the open physical index s. The four
    % explicit virtual wires carry the native directions. There are no
    % contractions: boundary signature (V,P)=(4,1), with two virtual
    % inputs and two virtual outputs.
    \begin{tenkzequation}
        \tenkzkernel
        \begin{tenkz}[rows={wire,wire,wire}, cols=3, bonds=none, frame=plane]
            \tn[at=(2,2), name=nativeSite, physical=up,
                ports={90:virtual,0:virtual,270:virtual,180:virtual},
                label pos=sw]{$a$}
            \tnwire[dir=from, name=nativeTop]{nativeSite.90}{open n}
            \tnwire[dir=to, name=nativeRight]{nativeSite.0}{open e}
            \tnwire[dir=to, name=nativeDown]{nativeSite.270}{open s}
            \tnwire[dir=from, name=nativeLeft]{nativeSite.180}{open w}
            \tnmark[form=label, label pos=n]{on nativeTop 1}{$\eta_t$}
            \tnmark[form=label, label pos=e]{on nativeRight 1}{$\eta_r$}
            \tnmark[form=label, label pos=s]{on nativeDown 1}{$\eta_b$}
            \tnmark[form=label, label pos=w]{on nativeLeft 1}{$\eta_l$}
            \tnmark[form=label, label pos=w]{on leg-n-2-2 1}{$s$}
        \end{tenkz}
    \end{tenkzequation}
    Then $\rho_v^o(g)_{\sigma,\eta}
        =\rho_\square(g)_{j_v\sigma,j_v\eta}$.
    For coefficient vectors, let $J_vx=x\circ j_v^{-1}$.
    The corresponding identities are
    $J_v\rho_v^o(g)=\rho_\square(g)J_v$ and
    $\mathcal A_v=\mathcal A J_v$.
    Hence $G$-injectivity of $\mathcal A$ for $\rho_\square$
    implies $G$-injectivity of every $\mathcal A_v$ for $\rho_v^o$.
    In numbered virtual coordinates the local map is
    $\mathcal A_vE_v$ and is $G$-injective for
    $E_v^{-1}\rho_v^oE_v$.
\end{theorem}
\begin{proof}\leanok
    \uses{thm:peps_native_torus_incidence}
    The outgoing indices $1,2$ contribute $U_{g^{-1}}^T$;
    indices $0,3$ contribute $U_g$. Reindexing the four incidence
    factors gives~(\ref{eq:peps_native_four_leg}). Reindexing the
    finite coefficient sum by $j_v$ gives both intertwining identities.
    Since $J_v$ is invertible, invariance and injectivity on invariant
    vectors pass from $\mathcal A$ to $\mathcal A_v$.
    Reindex once more by $E_v$ to obtain the numbered assertion.
\end{proof}

\begin{theorem}[A native $G$-injective averaging site exists]%
    \label{thm:peps_native_averaging_site}
    \lean{TNLean.PEPS.exists_isGInjective_torusSite}
    \leanok
    \uses{def:peps_g_injective}
    For every finite group representation $U$ on a finite alphabet
    $X$, there are a finite physical dimension $p$ and a native
    tensor $b:X^4\times\{0,\ldots,p-1\}\to\C$ which is
    $G$-injective for $\rho_\square$.
\end{theorem}
\begin{proof}\leanok
    Let $\Pi=|G|^{-1}\sum_g\rho_\square(g)$ and enumerate the
    coefficient basis of $\C^{X^4}$ by a linear isomorphism $E$
    to $\C^p$, with $p=|X|^4$. Take $b$ to be the matrix entries
    of $E\Pi$. The average satisfies $\Pi\rho_\square(g)=\Pi$
    and is the identity on invariant vectors. Composing with $E$
    therefore gives the required invariance and restricted injectivity.
\end{proof}

\subsection{Commuting closures satisfy the original plaquette constraints}

\begin{theorem}[Exterior insertions change only regional boundaries]%
    \label{thm:peps_native_inserted_region_support}
    \lean{TNLean.PEPS.graphInsertedBondNetwork_slice_mem_regionGroundSpace_of_internal_eq_one}
    \leanok
    \uses{def:peps_region_ground_space}
    On any finite simple graph, let $a_v$ be arbitrary site tensors
    with finite virtual alphabet $X$, and insert a matrix $K_e$ on
    each ordered bond. Suppose $K_e=1$ whenever both endpoints lie
    in $R$. For every fixed exterior physical configuration $\tau$,
    the resulting vector $\Psi_K(\,\cdot\,,\tau)$ on $R$ belongs
    to the original regional contraction range $\operatorname{im}\mathcal C_R$.
    No invertibility, group symmetry, or positivity of the inserted
    matrices is required.
\end{theorem}
\begin{proof}\leanok
    Split endpoint assignments into internal labels and pairs $\zeta$
    on all noninternal edges. The identity matrices on internal edges
    force their two endpoint labels to agree. For fixed $\zeta$, the
    remaining coefficient is
    \begin{align}
        c_\tau(\zeta)
         & =\prod_{e\notin E_R}K_e(\zeta_{e,h},\zeta_{e,t})
        \prod_{v\notin R}a_v(\zeta|_v,\tau_v),
        \label{eq:peps_native_inserted_exterior}            \\
        \Psi_K(\,\cdot\,,\tau)
         & =\sum_\zeta c_\tau(\zeta)
        \mathcal C_R(\theta(\zeta)).
        \label{eq:peps_native_inserted_slice}
    \end{align}
    Here $\theta(\zeta)$ selects the retained endpoint on every crossing
    edge. Each summand is a boundary column of the original regional
    map, so the sum belongs to its range.
\end{proof}

\begin{definition}[Native inserted matrices in ordered coordinates]%
    \label{def:peps_native_inserted_matrices}
    \lean{TNLean.PEPS.torusGraphBondMatrix,
        TNLean.PEPS.torusGraphBondMatrix_right,
        TNLean.PEPS.torusGraphBondMatrix_up}
    \leanok
    Let $O^h_v$ act on the native rightward bond from $v$, and
    $O^v_v$ on the downward bond from $(x,y+1)$ to $v=(x,y)$.
    For an ordered edge $e=(t,h)$ put
    \begin{align}
        K_{e_{\mathrm{right}}(v)} & =
        \begin{cases}O^h_v,&t=v,\\(O^h_v)^T,&t\ne v,\end{cases}
        \label{eq:peps_native_inserted_horizontal} \\
        K_{e_{\mathrm{up}}(v)}    & =
        \begin{cases}(O^v_v)^T,&t=v,\\O^v_v,&t\ne v.\end{cases}
        \label{eq:peps_native_inserted_vertical}
    \end{align}
    Thus a transpose is inserted exactly when the native arrow and
    ordered-edge arrow disagree.
\end{definition}

\begin{theorem}[Equality of native and inserted graph contractions]%
    \label{thm:peps_native_inserted_network}
    \lean{TNLean.PEPS.torusBondNetwork_eq_graphInsertedBondNetwork}
    \leanok
    \uses{def:peps_native_inserted_matrices,def:peps_native_incident_coordinates}
    For every physical configuration, the native torus contraction
    with matrices $O^h,O^v$ equals the ordered inserted graph
    contraction with the matrices $K$ of
    (\ref{eq:peps_native_inserted_horizontal})--
    (\ref{eq:peps_native_inserted_vertical}) and site tensors $a_v$.
\end{theorem}
\begin{proof}\leanok
    Reindex all native top/right/down/left labels by the incidence
    bijections $j_v$, then regroup them into ordered endpoint pairs.
    The horizontal and vertical bond coefficients become, respectively,
    \begin{align}
        O^h_v(\eta_{v+(1,0),\mathrm{left}},\eta_{v,\mathrm{right}}),
        \qquad
        O^v_v(\eta_{v,\mathrm{top}},\eta_{v+(0,1),\mathrm{down}}).
    \end{align}
    Equations~(\ref{eq:peps_native_inserted_horizontal}) and
    (\ref{eq:peps_native_inserted_vertical}) give exactly these entries
    in both possible vertex orders. The site factors are unchanged,
    so the reindexed finite sums agree.
\end{proof}

\begin{theorem}[Seams beyond a native plaquette]%
    \label{thm:peps_native_plaquette_seam_exclusion}
    \lean{TNLean.PEPS.torusGraphBondMatrix_closureAt_eq_one_on_plaquette_internal}
    \leanok
    \uses{def:peps_native_inserted_matrices}
    For $v=(x,y)$ let
    $R_v=v+\{(0,0),(1,0),(1,1),(0,1)\}$.
    Place the horizontal closure seam at $x+2$ and the vertical
    closure seam at $y+2$. For any $g,h\in G$, the corresponding
    inserted graph matrices are the identity on every edge internal
    to $R_v$. This assertion includes $m=3$ or $n=3$.
\end{theorem}
\begin{proof}\leanok
    The first coordinates in $R_v$ are $x,x+1$, whereas
    $x+2$ differs from both modulo $m\ge3$. The analogous statement
    holds for $y,y+1,y+2$ modulo $n$.
    An internal rightward edge cannot end at the horizontal seam,
    and an internal upward edge cannot end at the vertical seam.
    Their insertion matrices are therefore $1$; any transpose
    required by ordered coordinates leaves $1$ unchanged.
\end{proof}

\begin{theorem}[Invariant commuting closures are parent ground vectors]%
    \label{thm:peps_native_invariant_closure_membership}
    \lean{TNLean.PEPS.torusGClosure_slice_mem_plaquetteGroundSpace_of_invariant,
        TNLean.PEPS.torusGClosure_mem_torusParentGroundSpace_of_invariant,
        TNLean.PEPS.torusParent_annihilates_torusGClosure_of_invariant,
        TNLean.PEPS.torusParentHamiltonian_annihilates_torusGClosure_of_invariant,
        TNLean.PEPS.span_torusGClosure_commuting_le_torusParentKernel_of_invariant}
    \leanok
    \uses{def:peps_torus_g_closure,def:peps_region_parent_interaction}
    Let $U$ be any representation on finite $X$, and suppose
    $\mathcal A\rho_\square(k)=\mathcal A$ for every $k\in G$.
    If $gh=hg$, every physical slice of the native closure
    $\lvert a;(g,h)\rangle$ on $R_v$ lies in
    $\operatorname{im}\mathcal C_{R_v}$.
    Consequently the closure belongs to the full plaquette parent
    space. If $P_v\ge0$ and
    $\ker P_v=\operatorname{im}\mathcal C_{R_v}$, then every embedded
    $P_v$ annihilates it, and so does
    $H=\sum_v(P_v\otimes1_{R_v^c})$. In particular,
    \begin{align}
        \operatorname{span}\{\lvert a;(g,h)\rangle:gh=hg\}
        \subseteq\ker H.
        \label{eq:peps_native_closure_inclusion}
    \end{align}
    This implication uses only local virtual invariance, without
    semi-regularity or $G$-injectivity.
\end{theorem}
\begin{proof}\leanok
    \uses{thm:peps_torus_commuting_closure_seams,
        thm:peps_native_plaquette_seam_exclusion,
        thm:peps_native_inserted_network,
        thm:peps_native_inserted_region_support,
        thm:peps_region_parent_common_kernel}
    Local invariance and $gh=hg$ move both closure seams to $x+2,y+2$
    without changing the vector. Theorem~\ref{thm:peps_native_plaquette_seam_exclusion}
    makes every internal insertion on $R_v$ the identity.
    After the native-to-ordered contraction identity, formula
    (\ref{eq:peps_native_inserted_slice}) expresses each physical slice
    as a sum of unrestricted regional boundary columns.
    For the four-site plaquette, collecting equal boundary labels gives
    the contraction below. All four sites carry $a$.
    % Formula: closure slice = sum_theta c_tau(theta) C_{R_v}(theta),
    % with K_e=1 on every internal edge and arbitrary crossing/exterior K_e.
    % Source: SCP10 paper_v3.tex:1622--1647; native seam exclusion above.
    % Ink-to-index: top row y+1, bottom row y; columns x,x+1.
    % Unmarked internal bonds are identity contractions. Eight labelled
    % exterior virtual legs form theta; the four transverse physical legs
    % remain open. The patch has (V,P)=(8,4); summing theta against c_tau
    % gives the regional physical slice with (V,P)=(0,4).
    \begin{tenkzequation}
        \tenkzkernel
        $\displaystyle\lvert a;(g,h)\rangle(\,\cdot\,,\tau)
            =\sum_\theta c_\tau(\theta)$
        % Tenkz lattice dimensions are literal DSL tokens, not printed products.
        \begin{tenkz}[lattice={2x2}, frame=plane, physical=up] % chktex 29
            \tn[at=(1,1), name=nativePlaquette11,
                ports={180:virtual:$\theta_{l,1}$,0:virtual,
                        90:virtual:$\theta_{t,0}$,270:virtual}, label pos=sw]{}
            \tn[at=(1,2), name=nativePlaquette12,
                ports={180:virtual,0:virtual:$\theta_{r,1}$,
                        90:virtual:$\theta_{t,1}$,270:virtual}, label pos=sw]{}
            \tn[at=(2,1), name=nativePlaquette21,
                ports={180:virtual:$\theta_{l,0}$,0:virtual,
                        90:virtual,270:virtual:$\theta_{b,0}$}, label pos=sw]{}
            \tn[at=(2,2), name=nativePlaquette22,
                ports={180:virtual,0:virtual:$\theta_{r,0}$,
                        90:virtual,270:virtual:$\theta_{b,1}$}, label pos=sw]{}
        \end{tenkz}
    \end{tenkzequation}
    The coefficients $c_\tau$ include the exterior contraction and
    every crossing insertion. Each local term therefore annihilates
    every closure slice. Summing the terms gives $H\lvert a;(g,h)\rangle=0$;
    linearity proves~(\ref{eq:peps_native_closure_inclusion}).
\end{proof}

\subsection{Dimension and exhaustion of the full kernel}

\begin{theorem}[Full uniform semi-regular native parent dimension]%
    \label{thm:peps_native_uniform_semiregular_parent_dimension}
    \lean{TNLean.PEPS.IsGInjective.finrank_torusParentGroundSpace_of_isSemiRegular,
        TNLean.PEPS.IsGInjective.finrank_torusParentKernel_of_isSemiRegular}
    \leanok
    \uses{def:peps_semiregular,def:peps_g_injective,
        def:peps_native_incident_coordinates,def:peps_region_parent_interaction}
    Let $G$ be finite, let $U$ be one unitary semi-regular representation
    on finite $X$, and let $a$ be one $G$-injective native tensor for
    $\rho_\square$, repeated at every vertex of a torus with $m,n\ge3$.
    For each native plaquette $R_v$, let $P_v\ge0$ have kernel
    $\operatorname{im}\mathcal C_{R_v}$ and set
    $H=\sum_v(P_v\otimes1_{R_v^c})$.
    If $\mathcal K_G=\{(g,h)\in G^2:gh=hg\}/G$, with simultaneous
    conjugation, then
    \begin{align}
        \dim\mathcal P_{\{R_v\}}(a)=\dim\ker H=|\mathcal K_G|.
        \label{eq:peps_native_semiregular_dimension}
    \end{align}
    The block decomposition and its positive multiplicities are
    consequences of $U$; no expansion of a parent vector into
    closures is assumed. This is the stated uniform-torus case of
    \cite[Theorem~5.9]{Schuch2010PEPS}.
\end{theorem}
\begin{proof}\leanok
    \uses{thm:peps_native_oriented_g_injectivity,
        thm:peps_native_torus_incidence,
        thm:peps_oriented_semiregular_physical_parent,
        thm:peps_native_averaging_site,
        thm:peps_oriented_physical_canonical_parent,
        thm:peps_regular_torus_parent_ground_space,
        thm:peps_region_parent_common_kernel}
    Native incidence coordinates make every graph tensor
    $G$-injective for the actual rightward/downward orientation.
    The torus has degree four. Every vertex belongs to its plaquette,
    and every edge has both endpoints in a native plaquette.
    Theorem~\ref{thm:peps_oriented_semiregular_physical_parent}
    therefore identifies the entire parent space of $a$ with that
    of the oriented regular averaging tensor.
    Independently construct a native $G$-injective site $b$ for the
    regular representation using
    Theorem~\ref{thm:peps_native_averaging_site}.
    Fixed-representation comparison identifies its parent space with
    the same oriented regular canonical space. The already established
    regular native-torus theorem gives dimension $|\mathcal K_G|$
    for $b$. Thus
    \begin{align}
        \dim\mathcal P(a)
        =\dim\mathcal P(A^{L,1,o})
        =\dim\mathcal P(b)=|\mathcal K_G|.
        \label{eq:peps_native_dimension_comparison}
    \end{align}
    Finally, positivity and the genuine local kernel conditions give
    $\ker H=\mathcal P_{\{R_v\}}(a)$.
\end{proof}

\begin{theorem}[Full uniform semi-regular native parent spanning]%
    \label{thm:peps_native_uniform_semiregular_parent_spanning}
    \lean{TNLean.PEPS.IsGInjective.torusParentKernel_eq_commutingClosureSpan_of_isSemiRegular,
        TNLean.PEPS.IsGInjective.torusParentGroundSpace_eq_commutingClosureSpan_of_isSemiRegular}
    \leanok
    \uses{def:peps_torus_g_closure,def:peps_semiregular,
        def:peps_g_injective,def:peps_region_parent_interaction}
    Under the hypotheses of
    Theorem~\ref{thm:peps_native_uniform_semiregular_parent_dimension},
    \begin{align}
        \ker H=\mathcal P_{\{R_v\}}(a)
        =\operatorname{span}\{\lvert a;(g,h)\rangle:gh=hg\}.
        \label{eq:peps_native_semiregular_spanning}
    \end{align}
    Thus the actual commuting closure vectors exhaust the entire
    positive plaquette-parent kernel for one repeated $G$-injective
    tensor and one uniform unitary semi-regular representation on
    native tori of periods at least three.
    This is the corresponding case of
    \cite[Theorems~5.7 and~5.9]{Schuch2010PEPS}.
\end{theorem}
\begin{proof}\leanok
    \uses{thm:peps_native_invariant_closure_membership,
        thm:peps_native_uniform_semiregular_parent_dimension,
        thm:peps_g_injective_closure_dimension,
        thm:peps_region_parent_common_kernel}
    Let $\mathcal V$ be the span on the right of
    (\ref{eq:peps_native_semiregular_spanning}).
    Invariance gives $\mathcal V\subseteq\ker H$ by
    (\ref{eq:peps_native_closure_inclusion}).
    Closures in the same simultaneous-conjugacy class coincide,
    and semi-regular $G$-injectivity makes the vectors from distinct
    commuting classes linearly independent. Hence
    $\dim\mathcal V=|\mathcal K_G|$ by
    Theorem~\ref{thm:peps_g_injective_closure_dimension}.
    Equation~(\ref{eq:peps_native_semiregular_dimension}) gives
    $\dim\ker H=|\mathcal K_G|$. Inclusion and equal finite
    dimension imply $\mathcal V=\ker H$.
    Taking canonical positive plaquette interactions identifies the
    same span with the full intersection of regional range conditions.
\end{proof}
